\section{Contoh Terhitung}

\begin{example}[Serangan \textit{ciphertext-only} pada Caesar cipher]
Diketahui cipherteks $C = \code{KHOORZRUOG}$ tanpa diketahui kunci. Karena ruang kunci hanya 26 pergeseran, penyerang mencoba tiap kunci $k$ dengan aturan dekripsi
\[ P_i = (C_i - k) \bmod 26, \]
lalu memeriksa mana hasil yang bermakna. Hasil uji lengkapnya:
\begin{itemize}
    \item $k=0$: \code{KHOORZRUOG}; \quad $k=1$: \code{JGNNQYQTNF}
    \item $k=2$: \code{IFMMPXPSME}; \quad $k=3$: \code{HELLOWORLD}
    \item $k=4$: \code{GDKKNVNQKC}; \quad $k=5$: \code{FCJJMUMPJB}
    \item $k=6$: \code{EBIILTLOIA}; \quad $k=7$: \code{DAHHKSKNHZ}
    \item $k=8$: \code{CZGGJRJMGY}; \quad $k=9$: \code{BYFFIQILFX}
    \item $k=10$: \code{AXEEHPHKEW}; \quad $k=11$: \code{ZWDDGOGJDV}
    \item $k=12$: \code{YVCCFNFICU}; \quad $k=13$: \code{XUBBEMEHBT}
    \item $k=14$: \code{WTAADLDGAS}; \quad $k=15$: \code{VSZZCKCFZR}
    \item $k=16$: \code{URYYBJBEYQ}; \quad $k=17$: \code{TQXXAIADXP}
    \item $k=18$: \code{SPWWZHZCWO}; \quad $k=19$: \code{ROVVYGYBVN}
    \item $k=20$: \code{QNUUXFXAUM}; \quad $k=21$: \code{PMTTWEWZTL}
    \item $k=22$: \code{OLSSVDVYSK}; \quad $k=23$: \code{NKRRUCUXRJ}
    \item $k=24$: \code{MJQQTBTWQI}; \quad $k=25$: \code{LIPPSASVPH}
\end{itemize}
Hanya $k=3$ yang menghasilkan teks bermakna, yaitu \code{HELLOWORLD}. Seluruh proses selesai dalam
\[ t = \frac{26}{10^9} \approx 2{,}6 \times 10^{-8}\ \text{detik}, \]
jadi cipherteks-only berhasil karena ruang kuncinya terlalu kecil: penyerang memperoleh kunci sekaligus plainteks.
\end{example}

\begin{example}[Serangan \textit{chosen-plaintext} pada Vigen\`ere cipher]
Pada Vigen\`ere cipher berlaku $C_i = (P_i + K_i) \bmod 26$. Misalkan kunci yang sedang
dipakai adalah \code{SCRAM}, yang dalam bilangan menjadi
\[ K = (18,\ 2,\ 17,\ 0,\ 12). \]
Penyerang memilih plainteks $P = \code{AAAAA}$, yaitu $P = (0,0,0,0,0)$, dan mengirimkannya
ke mesin enkripsi. Karena penjumlahan dengan nol tidak mengubah apa pun,
\[ C_i = (0 + K_i) \bmod 26 = K_i, \]
sehingga sistem mengembalikan
\[ C = (18,\ 2,\ 17,\ 0,\ 12) = \code{SCRAM}. \]
Cipherteks yang keluar \emph{sama persis dengan kuncinya}. Satu kali query sudah cukup
untuk memulihkan seluruh kunci, sehingga seluruh pesan lain yang memakai kunci itu dapat
dibaca. Inilah yang membuat \textit{chosen-plaintext attack} jauh lebih berbahaya daripada
\textit{ciphertext-only attack}: penyerang tidak lagi menebak-nebak, ia menyusun pertanyaan
yang jawabannya adalah kunci itu sendiri \citep{katz2020}.
\end{example}

\begin{example}[Serangan \textit{timing} pada pemeriksaan kata sandi]
Sebuah fungsi pemeriksaan kata sandi membandingkan huruf demi huruf dan berhenti
segera setelah menemukan huruf yang berbeda. Untuk kata sandi sepanjang $n$ huruf dari
alfabet 26 huruf, pemeriksaan biasa menuntut sampai
\[ 26^n \]
percobaan. Untuk $n = 4$ diperoleh $26^4 = 456.976$ percobaan, dan untuk $n = 8$
angkanya melompat menjadi $26^8 = 208.827.064.576$ percobaan --- mustahil dilakukan.

Dengan mengukur waktu tanggap, penyerang dapat menyusuri huruf demi huruf. Huruf
pertama ditemukan setelah paling banyak 26 percobaan, lalu huruf kedua setelah paling
banyak 26 percobaan lagi, dan seterusnya, sehingga totalnya menjadi
\[ 26 \times n. \]
Untuk $n = 4$ hanya dibutuhkan $26 \times 4 = 104$ percobaan, dan untuk $n = 8$ hanya
$26 \times 8 = 208$ percobaan --- turun sekitar empat ribu kali lipat pada kasus
pertama. Perhatikan bahwa penyerang tidak memecahkan algoritma apa pun; ia hanya
memanfaatkan kenyataan bahwa lama eksekusi bergantung pada nilai rahasianya. Inilah
alasan mengapa fungsi pembanding kata sandi modern selalu ditulis dalam bentuk
\textbf{waktu tetap}, misalnya dengan meng-XOR seluruh huruf lalu memeriksa hasilnya
sekali saja di akhir \citep{stallings2017}.
\end{example}

\begin{example}[Kapan \textit{brute-force} menjadi tidak praktis]
Sebagai pembanding, cipher dengan kunci 56 bit memiliki $|\mathcal{K}| = 2^{56} \approx 7{,}21 \times 10^{16}$ kunci. Dengan kecepatan $10^9$ kunci/detik,
\[ t = \frac{2^{56}}{10^9} \approx 2{,}28\ \text{tahun}. \]
Bila kecepatan dinaikkan menjadi $10^{12}$ kunci/detik, waktu turun menjadi
\[ t = \frac{2^{56}}{10^{12}} \approx 7{,}21 \times 10^{4}\ \text{detik} \approx 20\ \text{jam}, \]
sehingga ruang kunci 56 bit tidak lagi memadai. Contoh ini sekaligus memperlihatkan
bahwa penilaian \textit{computationally secure} tidak dapat dilepaskan dari teknologi
yang tersedia pada zamannya: kunci yang dianggap mustahil dipecahkan pada tahun 1977
menjadi dapat disusuri dalam hitungan jam dua puluh tahun kemudian.
\end{example}

\begin{example}[Serangan \textit{replay} dan penangkalan dengan nomor urut]
Misalkan sebuah sistem perbankan menerima perintah transfer tanpa nomor urut. Nasabah
mengirim perintah ``transfer Rp1.000.000 ke rekening X'', dan penyerang menangkap
cipherteksnya. Karena sistem tidak menyimpan riwayat perintah, cipherteks itu dapat
dikirim ulang berkali-kali; setiap pengiriman ulang menghasilkan perintah yang sah
menurut sistem.

Bila perintah diberi nomor urut, penangkalannya menjadi sederhana. Sistem menyimpan
nomor urut terakhir yang diterima, misalkan 103, dan menolak setiap nomor yang tidak
lebih besar darinya. Serangan itu kini menjadi:
\begin{itemize}
    \item Perintah 101 diputar ulang setelah 103 diterima $\Rightarrow$ \textbf{ditolak},
          karena $101 \not> 103$.
    \item Perintah 104 dikirim secara sah $\Rightarrow$ \textbf{diterima}, nomor terakhir
          menjadi 104.
\end{itemize}
Perhatikan bahwa enkripsi tidak diubah sama sekali dan kuncinya tetap rahasia; yang
ditambahkan hanyalah sebuah mekanisme yang membuat tiap pesan tidak dapat diulang.
Inilah sebabnya integritas dan kesegaran pesan (\textit{freshness}) harus dirancang
sebagai bagian dari protokol, bukan diserahkan kepada cipher \citep{menezes1996}.
\end{example}
